% GENERATED FILE. Edit canonical lessons, then run npm run generate:books.
% canonical-source-hash: fnv1a64:e950840172cbe62d
% canonical-lessons: TE-C506-sahincu, TE-C506-tattuko, TE-C506-nigraham, TE-C506-manassanti, TE-C506-upasamanam

\chapter{Bearing It}
\label{ch:te-bearing-it}
\hlchaptermodality{506}

\begin{chapteropening}
\textbf{By the end of this chapter:} \emph{I can say that I put up with something, coped or found relief.}

Complete the last lesson of chapter 506: I can say that I put up with something, coped or found relief.
\end{chapteropening}

\section[\texorpdfstring{\te{సహించు} (sahiñcu)}{సహించు (sahiñcu)}]{\te{సహించు} (sahiñcu) — to tolerate; to put up with}

\label{lesson:TE-C506-sahincu}



\begin{quote}
Before the new one: say the Telugu for to stand by, then the Telugu for to console.
\end{quote}

\subsection*{You'll want to know: \te{సహించు}}
\textbf{\te{సహించు}} — \emph{sahiñcu} — \textquotedblleft{}to tolerate; to put up with\textquotedblright{}.

\begin{quote}
\te{అతని} \te{అబద్ధాలను} \te{నేను} \te{ఇక} \te{సహించను}.
\end{quote}

\begin{quote}
\emph{atani abaddhālanu nēnu ika sahiñcanu.}
\end{quote}

\begin{quote}
I will not put up with his lies any more.
\end{quote}

Past (I …): \te{సహించాను}, \emph{sahiñcānu}.

Often in the negative: \te{సహించను}, I will not stand for it. The noun \te{సహనం} is patience.

\begin{cousinweb}
\textbf{\te{సహించు}} (\emph{sahiñcu}) is built on Sanskrit \emph{sahatē}, \textquotedblleft{}bears, suffers\textquotedblright{} (CDIAL 13304), from the root \emph{sah}, with the Telugu verb ending -\te{ఇంచు}.
\end{cousinweb}

\begin{grammarlens}[title={What you've built}]
One more word for describing what you lived and felt.
\end{grammarlens}

\subsection*{Your turn}
\begin{itemize}
\raggedright
  \item \emph{Say it:} \emph{sahiñcu}
  \item \emph{Say it:} \emph{sahiñcu}, once more
  \item \emph{Say it:} \emph{atani abaddhālanu nēnu ika sahiñcanu.} — I will not put up with his lies any more.
  \item \emph{Recall:} say the Telugu for to stand by, then the Telugu for to console, then say \emph{sahiñcu} again
\end{itemize}

\begin{tcolorbox}[breakable,colback=teal!4,colframe=teal!35!black,title={Before you move on}]
What does \te{సహించు} mean? (\textquotedblleft{}To tolerate; to put up with\textquotedblright{}.) Say it once more.
\end{tcolorbox}
\section[\texorpdfstring{\te{తట్టుకో} (taṭṭukō)}{తట్టుకో (taṭṭukō)}]{\te{తట్టుకో} (taṭṭukō) — to withstand; to cope with}

\label{lesson:TE-C506-tattuko}



\begin{quote}
Before the new one: say the Telugu for to console, then the Telugu for to tolerate.
\end{quote}

\subsection*{You'll want to know: \te{తట్టుకో}}
\textbf{\te{తట్టుకో}} — \emph{taṭṭukō} — \textquotedblleft{}to withstand; to cope with\textquotedblright{}.

\begin{quote}
\te{ఆ} \te{చలిని} \te{ముసలివాళ్ళు} \te{తట్టుకోలేరు}.
\end{quote}

\begin{quote}
\emph{ā calini musalivāḷḷu taṭṭukōlēru.}
\end{quote}

\begin{quote}
Old people cannot withstand that cold.
\end{quote}

Past (I …): \te{తట్టుకున్నాను}, \emph{taṭṭukunnānu}.

Mostly with -\te{లే}-: \te{తట్టుకోలేను}, I cannot bear it.

\begin{cousinweb}
\textbf{\te{తట్టుకో}} (\emph{taṭṭukō}) is a Telugu verb ending in the reflexive -\te{కో}; the origin of its stem is not settled.
\end{cousinweb}

\begin{grammarlens}[title={What you've built}]
One more word for describing what you lived and felt.
\end{grammarlens}

\subsection*{Your turn}
\begin{itemize}
\raggedright
  \item \emph{Say it:} \emph{taṭṭukō}
  \item \emph{Say it:} \emph{taṭṭukō}, once more
  \item \emph{Say it:} \emph{ā calini musalivāḷḷu taṭṭukōlēru.} — Old people cannot withstand that cold.
  \item \emph{Recall:} say the Telugu for to console, then the Telugu for to tolerate, then say \emph{taṭṭukō} again
\end{itemize}

\begin{tcolorbox}[breakable,colback=teal!4,colframe=teal!35!black,title={Before you move on}]
What does \te{తట్టుకో} mean? (\textquotedblleft{}To withstand; to cope with\textquotedblright{}.) Say it once more.
\end{tcolorbox}
\section[\texorpdfstring{\te{నిగ్రహం} (nigrahaṁ)}{నిగ్రహం (nigrahaṁ)}]{\te{నిగ్రహం} (nigrahaṁ) — self-restraint; self-control}

\label{lesson:TE-C506-nigraham}



\begin{quote}
Before the new one: say the Telugu for to tolerate, then the Telugu for to withstand.
\end{quote}

\subsection*{You'll want to know: \te{నిగ్రహం}}
\textbf{\te{నిగ్రహం}} — \emph{nigrahaṁ} — \textquotedblleft{}self-restraint; self-control\textquotedblright{}.

\begin{quote}
\te{కోపం} \te{వచ్చినా} \te{ఆమె} \te{నిగ్రహంతో} \te{మాట్లాడింది}.
\end{quote}

\begin{quote}
\emph{kōpaṁ vaccinā āme nigrahantō māṭlāḍindi.}
\end{quote}

\begin{quote}
Though she was angry, she spoke with restraint.
\end{quote}

Holding back anger or a craving: \te{నిగ్రహంతో}, with restraint. A formal word.

\begin{cousinweb}
\textbf{\te{నిగ్రహం}} (\emph{nigrahaṁ}) is borrowed from Sanskrit \emph{nigraha}, \textquotedblleft{}restraint; a rebuke\textquotedblright{} (CDIAL 7169), from \emph{ni-grah}, \textquotedblleft{}to hold down\textquotedblright{}: holding yourself down.
\end{cousinweb}

\begin{grammarlens}[title={What you've built}]
One more word for describing what you lived and felt.
\end{grammarlens}

\subsection*{Your turn}
\begin{itemize}
\raggedright
  \item \emph{Say it:} \emph{nigrahaṁ}
  \item \emph{Say it:} \emph{nigrahaṁ}, once more
  \item \emph{Say it:} \emph{kōpaṁ vaccinā āme nigrahantō māṭlāḍindi.} — Though she was angry, she spoke with restraint.
  \item \emph{Recall:} say the Telugu for to tolerate, then the Telugu for to withstand, then say \emph{nigrahaṁ} again
\end{itemize}

\begin{tcolorbox}[breakable,colback=teal!4,colframe=teal!35!black,title={Before you move on}]
What does \te{నిగ్రహం} mean? (\textquotedblleft{}Self-restraint; self-control\textquotedblright{}.) Say it once more.
\end{tcolorbox}
\section[\texorpdfstring{\te{మనశ్శాంతి} (manaśśānti)}{మనశ్శాంతి (manaśśānti)}]{\te{మనశ్శాంతి} (manaśśānti) — peace of mind}

\label{lesson:TE-C506-manassanti}



\begin{quote}
Before the new one: say the Telugu for to withstand, then the Telugu for self-restraint.
\end{quote}

\subsection*{You'll want to know: \te{మనశ్శాంతి}}
\textbf{\te{మనశ్శాంతి}} — \emph{manaśśānti} — \textquotedblleft{}peace of mind\textquotedblright{}.

\begin{quote}
\te{పల్లెలో} \te{ఉంటే} \te{నాకు} \te{మనశ్శాంతిగా} \te{ఉంటుంది}.
\end{quote}

\begin{quote}
\emph{pallelō uṇṭē nāku manaśśāntigā uṇṭundi.}
\end{quote}

\begin{quote}
In the village I have peace of mind.
\end{quote}

\te{మనసు}, mind, + \te{శాంతి}, peace.

\begin{cousinweb}
\textbf{\te{మనశ్శాంతి}} (\emph{manaśśānti}) is borrowed from Sanskrit: \emph{manas}, \textquotedblleft{}mind\textquotedblright{}, and \emph{śānti}, \textquotedblleft{}rest, peace\textquotedblright{} (CDIAL 12391); the \emph{s} of \emph{manas} becomes \emph{ś} before \emph{ś}.
\end{cousinweb}

\begin{grammarlens}[title={What you've built}]
One more word for describing what you lived and felt.
\end{grammarlens}

\subsection*{Your turn}
\begin{itemize}
\raggedright
  \item \emph{Say it:} \emph{manaśśānti}
  \item \emph{Say it:} \emph{manaśśānti}, once more
  \item \emph{Say it:} \emph{pallelō uṇṭē nāku manaśśāntigā uṇṭundi.} — In the village I have peace of mind.
  \item \emph{Recall:} say the Telugu for to withstand, then the Telugu for self-restraint, then say \emph{manaśśānti} again
\end{itemize}

\begin{tcolorbox}[breakable,colback=teal!4,colframe=teal!35!black,title={Before you move on}]
What does \te{మనశ్శాంతి} mean? (\textquotedblleft{}Peace of mind\textquotedblright{}.) Say it once more.
\end{tcolorbox}
\section[\texorpdfstring{\te{ఉపశమనం} (upaśamanaṁ)}{ఉపశమనం (upaśamanaṁ)}]{\te{ఉపశమనం} (upaśamanaṁ) — relief; respite}

\label{lesson:TE-C506-upasamanam}



\begin{quote}
Before the new one: say the Telugu for self-restraint, then the Telugu for peace of mind.
\end{quote}

\subsection*{You'll want to know: \te{ఉపశమనం}}
\textbf{\te{ఉపశమనం}} — \emph{upaśamanaṁ} — \textquotedblleft{}relief; respite\textquotedblright{}.

\begin{quote}
\te{మందు} \te{వేసుకున్నాక} \te{నొప్పికి} \te{కొంచెం} \te{ఉపశమనం} \te{దొరికింది}.
\end{quote}

\begin{quote}
\emph{mandu vēsukunnāka noppiki koñceṁ upaśamanaṁ dorikindi.}
\end{quote}

\begin{quote}
After I took the medicine, the pain eased a little.
\end{quote}

\te{ఉపశమనం} \te{కలుగు} or \te{దొరుకు}, relief comes.

\begin{cousinweb}
\textbf{\te{ఉపశమనం}} (\emph{upaśamanaṁ}) is borrowed from Sanskrit \emph{upaśamana}, \textquotedblleft{}calming, appeasing, mitigation\textquotedblright{} (Monier-Williams).
\end{cousinweb}

\begin{grammarlens}[title={What you've built}]
One more word for describing what you lived and felt.
\end{grammarlens}

\subsection*{Your turn}
\begin{itemize}
\raggedright
  \item \emph{Say it:} \emph{upaśamanaṁ}
  \item \emph{Say it:} \emph{upaśamanaṁ}, once more
  \item \emph{Say it:} \emph{mandu vēsukunnāka noppiki koñceṁ upaśamanaṁ dorikindi.} — After I took the medicine, the pain eased a little.
  \item \emph{Recall:} say the Telugu for self-restraint, then the Telugu for peace of mind, then say \emph{upaśamanaṁ} again
\end{itemize}

\begin{tcolorbox}[breakable,colback=teal!4,colframe=teal!35!black,title={Before you move on}]
What does \te{ఉపశమనం} mean? (\textquotedblleft{}Relief; respite\textquotedblright{}.) Say it once more.
\end{tcolorbox}
